\documentclass{situationist-preprint}

\title{Memory Is Not Necessary}
\msid{PID-08-b}
\msfamily{PID-08}
\msdiscipline{Personal Identity and the Self}
\mstype{Companion manuscript}
\mskeywords{amnesia; memory criterion; self-concern; Williams; Reid; personal identity}

\begin{document}
\maketitle

\begin{abstract}
PID-08-a argues that a person can survive total amnesia, from the defects of the memory criterion and from clinical cases. We add an argument from anticipation. A person told that she will lose all her autobiographical memory and then suffer pain fears the pain, and her fear is not irrational; Williams made the point to show how hard it is to separate oneself from one's body. We argue that the fear is warranted by the psychological criterion of PID-01 without memory, and that a gradual version of the case leaves the memory theorist no non-arbitrary point at which the person ceases to exist.
\end{abstract}

\section{The argument from anticipation}

\citet{Williams1970} asked us to imagine being told that tomorrow we will be tortured, and that before the torture we will lose all our memories. The prospect of amnesia does not make the prospect of torture less frightening; if anything it adds to the fear. Williams concluded that our anticipation of future experience is not governed by memory.

\begin{proposition}\label{p:fear}
If memory continuity were necessary for survival, a person told that she will undergo total amnesia and then pain would have no self-interested reason to fear the pain. She has such a reason. So memory continuity is not necessary for survival.
\end{proposition}

\begin{proof}
Self-interested fear of a future pain is warranted only if the one who will suffer it is oneself or stands to oneself in a relation that carries what matters. If memory were necessary for survival, the sufferer would not be she. Nor, on the memory theorist's own view, would anything else connect them that carries what matters, since the memory theorist holds that memory is what makes a future person's experiences one's own concern. But the sufferer has her character, values, desires, skills, knowledge and emotional dispositions, and her capacity for self-conscious thought (PID-08-a, Definition~2.2), realized in the same brain. Her fear of the pain is the fear of a person who knows that what will be hurt is the one who loves what she loves, will want what she wants and will feel the pain as she would. That fear is warranted, and it is warranted because the person who suffers will be she, by the psychological criterion of PID-01 with the connections of PID-08-a, Assumption~2.3.
\end{proof}

\section{The gradual case}

Suppose memories are removed one at a time over a year, the other capacities continuing. The memory theorist must say that at some point the person ceases to exist and another begins, since at the end no memory connects them. Yet at no step is anything removed but a single memory, and nothing else changes. Either some single memory makes the difference between survival and death, which is incredible, or memory continuity is not necessary. \citet[III.6]{Reid1785} pressed against Locke the case of the brave officer, which shows that memory connections are not transitive while identity is; the gradual case shows, in addition, that the absence of all memory connections at the end of a sequence of trivial changes cannot mark a death.

\begin{objection}
In Williams's case the fear may be fear for the body, not for oneself. It shows only that people are attached to their bodies.
\end{objection}

\begin{reply}
If the fear were for the body, it would be equally apt if the body were to be damaged after the whole of one's psychology had been irreversibly destroyed, as in a permanent vegetative state. It is not: one may regret such damage, but one does not anticipate it as pain one will feel. What the anticipation tracks is the continuation of the psychology, minus the memories, in whatever realizes it.
\end{reply}

\section{Conclusion}

The rational fear of pain after amnesia, and the absence of any point in gradual amnesia at which a person could be said to die, both show that memory continuity is not necessary for survival. The other psychological capacities suffice.

\begin{thebibliography}{99}\small
\bibitem[Reid(1785)]{Reid1785} Reid, T. (1785). \emph{Essays on the Intellectual Powers of Man}. Cited by essay and chapter.
\bibitem[Williams(1970)]{Williams1970} Williams, B. (1970). The self and the future. \emph{The Philosophical Review}, 79(2), 161--180.
\end{thebibliography}

\end{document}
